\documentclass[11pt,a4paper]{article}
\usepackage[utf8]{inputenc}
\usepackage[T1]{fontenc}
\usepackage{amsmath,amssymb}
\usepackage{graphicx}
\usepackage{booktabs}
\usepackage{siunitx}
\usepackage{geometry}
\geometry{margin=1in}
\usepackage{pythontex}
\usepackage{hyperref}
\usepackage{float}
\usepackage{cite}

\title{Reaction Engineering\\CSTR and PFR Design}
\author{Chemical Engineering Research Group}
\date{\today}

\begin{document}
\maketitle

\begin{abstract}
This report presents comprehensive computational analysis of chemical reactor design and kinetics, focusing on batch reactors, continuous stirred tank reactors (CSTRs), and plug flow reactors (PFRs). We investigate first-order and second-order reaction kinetics, residence time distributions, temperature effects via Arrhenius behavior, and reactor performance comparisons using Levenspiel plots. The analysis provides design equations, numerical solutions for non-isothermal operation, and optimization strategies for industrial reactor systems.
\end{abstract}

\section{Introduction}

Chemical reactor design is fundamental to the chemical process industry, determining the size, configuration, and operating conditions required to achieve desired conversion and selectivity \cite{Fogler2016, Levenspiel1999}. The choice between batch, CSTR, and PFR configurations depends on reaction kinetics, heat transfer requirements, and production scale \cite{Smith2005}.

This computational study examines reactor performance across different configurations, investigating the effects of reaction order, temperature, and residence time on conversion. We employ Python-based numerical methods to solve design equations and generate Levenspiel plots for reactor comparison \cite{Scott2006}.

\begin{pycode}

import numpy as np
import matplotlib.pyplot as plt
from scipy.integrate import odeint, solve_ivp
from scipy.optimize import fsolve, brentq
plt.rcParams['text.usetex'] = True
plt.rcParams['font.family'] = 'serif'

# Global constants
R = 8.314  # J/(mol·K) - Universal gas constant

\end{pycode}

\section{Batch Reactor Kinetics}

For a constant-volume batch reactor, the design equation for species A is:
\begin{equation}
\frac{dC_A}{dt} = -r_A
\end{equation}

For a first-order irreversible reaction $A \rightarrow P$ with rate constant $k$:
\begin{equation}
-r_A = k C_A \quad \Rightarrow \quad C_A(t) = C_{A0} e^{-kt}
\end{equation}

For a second-order reaction $A \rightarrow P$:
\begin{equation}
-r_A = k C_A^2 \quad \Rightarrow \quad C_A(t) = \frac{C_{A0}}{1 + k C_{A0} t}
\end{equation}

The conversion $X_A$ is defined as: $X_A = 1 - C_A/C_{A0}$

\begin{pycode}
# Batch reactor kinetics comparison
k1 = 0.1  # rate constant for first-order (1/min)
k2 = 0.15  # rate constant for second-order (L/(mol·min))
C_A0 = 1.0  # initial concentration (mol/L)
t = np.linspace(0, 60, 500)

# First-order kinetics
C_A_first = C_A0 * np.exp(-k1 * t)
X_A_first = 1 - C_A_first / C_A0

# Second-order kinetics
C_A_second = C_A0 / (1 + k2 * C_A0 * t)
X_A_second = 1 - C_A_second / C_A0

# Store key results
t_90_first = -np.log(0.1) / k1
t_90_second = 9 / (k2 * C_A0)

fig, (ax1, ax2) = plt.subplots(1, 2, figsize=(12, 5))

# Concentration profiles
ax1.plot(t, C_A_first, 'b-', linewidth=2, label='First-order')
ax1.plot(t, C_A_second, 'r--', linewidth=2, label='Second-order')
ax1.axhline(y=0.1*C_A0, color='gray', linestyle=':', alpha=0.5)
ax1.set_xlabel('Time (min)', fontsize=11)
ax1.set_ylabel('$C_A$ (mol/L)', fontsize=11)
ax1.set_title('Batch Reactor Concentration Decay', fontsize=12)
ax1.legend(fontsize=10)
ax1.grid(True, alpha=0.3)

# Conversion profiles
ax2.plot(t, X_A_first, 'b-', linewidth=2, label='First-order')
ax2.plot(t, X_A_second, 'r--', linewidth=2, label='Second-order')
ax2.axhline(y=0.9, color='gray', linestyle=':', alpha=0.5, label='90\\% conversion')
ax2.set_xlabel('Time (min)', fontsize=11)
ax2.set_ylabel('Conversion $X_A$', fontsize=11)
ax2.set_title('Batch Reactor Conversion', fontsize=12)
ax2.legend(fontsize=10)
ax2.grid(True, alpha=0.3)

plt.tight_layout()
plt.savefig('reaction_engineering_plot1.pdf', dpi=150, bbox_inches='tight')
plt.close()
\end{pycode}

\begin{figure}[H]
\centering
\includegraphics[width=0.95\textwidth]{reaction_engineering_plot1.pdf}
\caption{Batch reactor performance for first-order and second-order reactions. First-order kinetics exhibit exponential decay, while second-order reactions show hyperbolic behavior. For 90\% conversion with $C_{A0} = \py{C_A0}$ mol/L, first-order requires $\py{f'{t_90_first:.1f}'}$ min while second-order requires $\py{f'{t_90_second:.1f}'}$ min, demonstrating the impact of reaction order on batch time.}
\end{figure}

\section{CSTR Design}

For a continuous stirred tank reactor (CSTR) at steady state, the design equation is:
\begin{equation}
\tau = \frac{V}{Q} = \frac{C_{A0} - C_A}{-r_A} = \frac{C_{A0} X_A}{-r_A}
\end{equation}

For a first-order reaction in a single CSTR:
\begin{equation}
\tau = \frac{X_A}{k(1 - X_A)}
\end{equation}

For $N$ equal-volume CSTRs in series, each with volume $V_i = V/N$:
\begin{equation}
X_{A,N} = 1 - \frac{1}{(1 + k\tau/N)^N}
\end{equation}

\begin{pycode}
# CSTR design - single and multiple reactors in series
k_cstr = 0.1  # rate constant (1/min)
X_target = 0.90  # target conversion

# Single CSTR residence time for target conversion
tau_single = X_target / (k_cstr * (1 - X_target))

# Multiple CSTRs in series
N_reactors = np.array([1, 2, 3, 4, 5, 10])
X_A_range = np.linspace(0.01, 0.99, 200)

fig, (ax1, ax2) = plt.subplots(1, 2, figsize=(12, 5))

# Conversion vs residence time for different numbers of CSTRs
tau_values = np.linspace(0, 50, 300)
for N in [1, 2, 3, 5]:
    if N == 1:
        X_vals = (k_cstr * tau_values) / (1 + k_cstr * tau_values)
        label_str = 'Single CSTR'
    else:
        X_vals = 1 - 1 / (1 + k_cstr * tau_values / N)**N
        label_str = f'{N} CSTRs in series'
    ax1.plot(tau_values, X_vals, linewidth=2, label=label_str)

ax1.axhline(y=0.9, color='gray', linestyle=':', alpha=0.5)
ax1.axvline(x=tau_single, color='gray', linestyle=':', alpha=0.5)
ax1.set_xlabel('Total Residence Time $\\tau$ (min)', fontsize=11)
ax1.set_ylabel('Conversion $X_A$', fontsize=11)
ax1.set_title('CSTR Performance: Effect of Staging', fontsize=12)
ax1.legend(fontsize=9)
ax1.grid(True, alpha=0.3)
ax1.set_xlim(0, 50)

# Volume requirement for different staging configurations
total_volumes = []
for N in N_reactors:
    if N == 1:
        tau_req = X_target / (k_cstr * (1 - X_target))
    else:
        # Solve for tau: X_A = 1 - 1/(1 + k*tau/N)^N
        def equation(tau):
            return 1 - 1/(1 + k_cstr*tau/N)**N - X_target
        tau_req = brentq(equation, 0, 100)
    total_volumes.append(tau_req)  # Assuming Q=1 so V=tau

ax2.plot(N_reactors, total_volumes, 'o-', linewidth=2, markersize=8, color='steelblue')
ax2.axhline(y=tau_single, color='gray', linestyle='--', alpha=0.5, label='Single CSTR')
ax2.set_xlabel('Number of CSTRs in Series', fontsize=11)
ax2.set_ylabel('Total Volume Required (L)', fontsize=11)
ax2.set_title(f'Volume Optimization for $X_A = {X_target}$', fontsize=12)
ax2.legend(fontsize=10)
ax2.grid(True, alpha=0.3)

plt.tight_layout()
plt.savefig('reaction_engineering_plot2.pdf', dpi=150, bbox_inches='tight')
plt.close()

# Store results
V_single = total_volumes[0]
V_three = total_volumes[2]
reduction_pct = (V_single - V_three) / V_single * 100
\end{pycode}

\begin{figure}[H]
\centering
\includegraphics[width=0.95\textwidth]{reaction_engineering_plot2.pdf}
\caption{CSTR staging analysis showing the benefit of multiple reactors in series. For first-order kinetics with $k = \py{k_cstr}$ min$^{-1}$, achieving $X_A = \py{X_target}$ requires $\tau = \py{f'{tau_single:.1f}'}$ min for a single CSTR. Using three CSTRs in series reduces total volume by $\py{f'{reduction_pct:.1f}'}$\%, approaching PFR performance while maintaining the advantages of perfect mixing within each stage.}
\end{figure}

\section{Plug Flow Reactor (PFR) Design}

For a PFR, the differential design equation is:
\begin{equation}
\frac{dX_A}{dV} = \frac{-r_A}{F_{A0}}
\end{equation}

For isothermal first-order kinetics, this integrates to:
\begin{equation}
V = \frac{F_{A0}}{k} \ln\left(\frac{1}{1-X_A}\right) = \frac{Q C_{A0}}{k} \ln\left(\frac{1}{1-X_A}\right)
\end{equation}

Comparing residence times: $\tau_{PFR} = -\frac{1}{k}\ln(1-X_A)$ vs. $\tau_{CSTR} = \frac{X_A}{k(1-X_A)}$

For non-isothermal operation, we must solve coupled ODEs numerically:
\begin{align}
\frac{dX_A}{dV} &= \frac{k(T) C_{A0} (1-X_A)}{F_{A0}} \\
\frac{dT}{dV} &= \frac{(-\Delta H_R) k(T) C_{A0} (1-X_A) - UA(T - T_c)}{\rho C_p Q}
\end{align}

\begin{pycode}
# PFR design with non-isothermal operation
k_pfr = 0.1  # rate constant at reference T (1/min)
C_A0_pfr = 1.0  # inlet concentration (mol/L)
Q_pfr = 10.0  # volumetric flow rate (L/min)
F_A0 = Q_pfr * C_A0_pfr  # molar flow rate (mol/min)

# Isothermal PFR
X_A_pfr = np.linspace(0.01, 0.95, 100)
V_pfr = (F_A0 / k_pfr) * np.log(1 / (1 - X_A_pfr))
tau_pfr = V_pfr / Q_pfr

# Compare with CSTR
tau_cstr_comp = X_A_pfr / (k_pfr * (1 - X_A_pfr))
V_cstr_comp = tau_cstr_comp * Q_pfr

# Non-isothermal PFR with exothermic reaction
def pfr_nonisothermal(V, y):
    X_A, T = y

    # Arrhenius temperature dependence
    A = 1e10  # pre-exponential factor (1/min)
    Ea = 75000  # activation energy (J/mol)
    k_T = A * np.exp(-Ea / (R * T))

    # Energy balance parameters
    DH_R = -50000  # heat of reaction (J/mol) - exothermic
    rho = 1000  # density (kg/m³)
    Cp = 4180  # heat capacity (J/(kg·K))
    U = 500  # heat transfer coefficient (W/(m²·K))
    A_heat = 0.5  # heat transfer area per volume (m²/L)
    T_c = 298  # coolant temperature (K)

    dX_dV = (k_T * C_A0_pfr * (1 - X_A)) / F_A0
    dT_dV = ((-DH_R) * k_T * C_A0_pfr * (1 - X_A) - U * A_heat * (T - T_c)) / (rho * Cp * Q_pfr / 1000)

    return [dX_dV, dT_dV]

# Initial conditions
T_inlet = 320  # inlet temperature (K)
V_range = np.linspace(0, 50, 500)

sol = odeint(pfr_nonisothermal, [0, T_inlet], V_range)
X_A_noniso = sol[:, 0]
T_profile = sol[:, 1]

fig, (ax1, ax2) = plt.subplots(1, 2, figsize=(12, 5))

# Isothermal reactor comparison
ax1.plot(tau_pfr, X_A_pfr, 'b-', linewidth=2, label='PFR')
ax1.plot(tau_cstr_comp, X_A_pfr, 'r--', linewidth=2, label='CSTR')
ax1.axhline(y=0.9, color='gray', linestyle=':', alpha=0.5)
ax1.set_xlabel('Residence Time $\\tau$ (min)', fontsize=11)
ax1.set_ylabel('Conversion $X_A$', fontsize=11)
ax1.set_title('PFR vs CSTR: Isothermal First-Order', fontsize=12)
ax1.legend(fontsize=10)
ax1.grid(True, alpha=0.3)
ax1.set_xlim(0, 50)

# Non-isothermal PFR
ax2_temp = ax2.twinx()
ax2.plot(V_range, X_A_noniso, 'b-', linewidth=2, label='Conversion')
ax2_temp.plot(V_range, T_profile, 'r--', linewidth=2, label='Temperature')
ax2.set_xlabel('Reactor Volume (L)', fontsize=11)
ax2.set_ylabel('Conversion $X_A$', fontsize=11, color='b')
ax2_temp.set_ylabel('Temperature (K)', fontsize=11, color='r')
ax2.set_title('Non-Isothermal PFR with Cooling', fontsize=12)
ax2.tick_params(axis='y', labelcolor='b')
ax2_temp.tick_params(axis='y', labelcolor='r')
ax2.grid(True, alpha=0.3)

# Combine legends
lines1, labels1 = ax2.get_legend_handles_labels()
lines2, labels2 = ax2_temp.get_legend_handles_labels()
ax2.legend(lines1 + lines2, labels1 + labels2, fontsize=10, loc='center right')

plt.tight_layout()
plt.savefig('reaction_engineering_plot3.pdf', dpi=150, bbox_inches='tight')
plt.close()

# Calculate volumes for 90% conversion
V_pfr_90 = (F_A0 / k_pfr) * np.log(1 / 0.1)
V_cstr_90 = (0.9 / (k_pfr * 0.1)) * Q_pfr
volume_ratio = V_cstr_90 / V_pfr_90
T_max = np.max(T_profile)
T_rise = T_max - T_inlet
\end{pycode}

\begin{figure}[H]
\centering
\includegraphics[width=0.95\textwidth]{reaction_engineering_plot3.pdf}
\caption{PFR analysis showing superior performance over CSTR for positive-order reactions. For 90\% conversion, PFR requires $V = \py{f'{V_pfr_90:.1f}'}$ L while CSTR needs $\py{f'{V_cstr_90:.1f}'}$ L, a $\py{f'{volume_ratio:.1f}'}$× difference. The non-isothermal case demonstrates temperature rise of $\py{f'{T_rise:.1f}'}$ K due to exothermic reaction ($\Delta H_R = -50$ kJ/mol), requiring active cooling to prevent thermal runaway.}
\end{figure}

\section{Arrhenius Temperature Dependence}

The Arrhenius equation describes the temperature dependence of the rate constant:
\begin{equation}
k(T) = A \exp\left(-\frac{E_a}{RT}\right)
\end{equation}

where $A$ is the pre-exponential factor, $E_a$ is the activation energy, $R$ is the gas constant, and $T$ is absolute temperature.

Taking the logarithm yields a linear form:
\begin{equation}
\ln k = \ln A - \frac{E_a}{R} \cdot \frac{1}{T}
\end{equation}

This enables determination of $E_a$ from the slope of an Arrhenius plot ($\ln k$ vs. $1/T$).

\begin{pycode}
# Arrhenius temperature dependence analysis
A_arr = 1e10  # pre-exponential factor (1/min)
Ea_arr = 75000  # activation energy (J/mol)
R_const = 8.314  # J/(mol·K)

# Temperature range
T_range_C = np.linspace(0, 100, 200)
T_range_K = T_range_C + 273.15

# Calculate rate constants
k_values = A_arr * np.exp(-Ea_arr / (R_const * T_range_K))

# Reference temperatures
T_ref = 298.15  # 25°C
T_high = 373.15  # 100°C
k_ref = A_arr * np.exp(-Ea_arr / (R_const * T_ref))
k_high = A_arr * np.exp(-Ea_arr / (R_const * T_high))
k_ratio = k_high / k_ref

# Arrhenius plot data for parameter estimation
T_expt = np.array([273, 298, 323, 348, 373])  # K
k_expt = A_arr * np.exp(-Ea_arr / (R_const * T_expt))
# Add some noise to simulate experimental data
np.random.seed(42)
k_expt_noisy = k_expt * (1 + 0.05 * np.random.randn(len(k_expt)))

# Linear regression for Ea determination
inv_T = 1 / T_expt
ln_k = np.log(k_expt_noisy)
coeffs = np.polyfit(inv_T, ln_k, 1)
Ea_fitted = -coeffs[0] * R_const
A_fitted = np.exp(coeffs[1])
ln_k_fit = np.polyval(coeffs, inv_T)

fig, (ax1, ax2) = plt.subplots(1, 2, figsize=(12, 5))

# Rate constant vs temperature
ax1.semilogy(T_range_C, k_values, 'b-', linewidth=2)
ax1.plot([T_ref-273.15, T_high-273.15], [k_ref, k_high], 'ro', markersize=8, label='Reference points')
ax1.set_xlabel('Temperature (°C)', fontsize=11)
ax1.set_ylabel('Rate Constant $k$ (min$^{-1}$)', fontsize=11)
ax1.set_title('Arrhenius Temperature Dependence', fontsize=12)
ax1.legend(fontsize=10)
ax1.grid(True, alpha=0.3, which='both')

# Arrhenius plot for parameter estimation
ax2.plot(inv_T * 1000, ln_k, 'o', markersize=8, label='Experimental data')
ax2.plot(inv_T * 1000, ln_k_fit, 'r--', linewidth=2, label=f'Fit: $E_a$ = {Ea_fitted/1000:.1f} kJ/mol')
ax2.set_xlabel('1000/T (K$^{-1}$)', fontsize=11)
ax2.set_ylabel('$\\ln(k)$', fontsize=11)
ax2.set_title('Arrhenius Plot for Parameter Estimation', fontsize=12)
ax2.legend(fontsize=10)
ax2.grid(True, alpha=0.3)

plt.tight_layout()
plt.savefig('reaction_engineering_plot4.pdf', dpi=150, bbox_inches='tight')
plt.close()

# Calculate reaction time sensitivity
X_target_arr = 0.90
t_25C = -np.log(1 - X_target_arr) / k_ref
t_50C = -np.log(1 - X_target_arr) / (A_arr * np.exp(-Ea_arr / (R_const * 323.15)))
time_reduction = (t_25C - t_50C) / t_25C * 100
\end{pycode}

\begin{figure}[H]
\centering
\includegraphics[width=0.95\textwidth]{reaction_engineering_plot4.pdf}
\caption{Arrhenius analysis showing exponential temperature dependence of reaction rate. With $E_a = \py{f'{Ea_arr/1000:.0f}'}$ kJ/mol, the rate constant increases $\py{f'{k_ratio:.1f}'}$× from 25°C to 100°C. For batch operation achieving 90\% conversion, increasing temperature from 25°C to 50°C reduces reaction time by $\py{f'{time_reduction:.1f}'}$\%, demonstrating the strong coupling between kinetics and thermal management in reactor design.}
\end{figure}

\section{Levenspiel Plots for Reactor Comparison}

The Levenspiel plot displays $1/(-r_A)$ versus conversion $X_A$, enabling graphical reactor design and comparison \cite{Levenspiel1999}. The area under the curve represents the required reactor volume:

For PFR: $V = F_{A0} \int_0^{X_A} \frac{dX_A}{-r_A}$ (area under curve)

For CSTR: $V = F_{A0} \frac{X_A}{-r_A}$ (area of rectangle)

This graphical method reveals that PFRs always require less volume than CSTRs for positive-order reactions.

\begin{pycode}
# Levenspiel plots for different reaction orders
F_A0_lev = 10.0  # molar flow rate (mol/min)
C_A0_lev = 1.0  # inlet concentration (mol/L)
X_vals = np.linspace(0.01, 0.95, 200)

# First-order reaction
k1_lev = 0.1  # 1/min
C_A_first = C_A0_lev * (1 - X_vals)
r_A_first = -k1_lev * C_A_first
inv_r_A_first = -1 / r_A_first

# Second-order reaction
k2_lev = 0.15  # L/(mol·min)
C_A_second = C_A0_lev * (1 - X_vals)
r_A_second = -k2_lev * C_A_second**2
inv_r_A_second = -1 / r_A_second

# Zero-order reaction
k0_lev = 0.1  # mol/(L·min)
r_A_zero = -k0_lev
inv_r_A_zero = -1 / r_A_zero * np.ones_like(X_vals)

# Calculate volumes for X = 0.80
X_design = 0.80
idx_design = np.argmin(np.abs(X_vals - X_design))

# PFR volumes (integration)
V_pfr_first = F_A0_lev * np.trapz(inv_r_A_first[:idx_design], X_vals[:idx_design])
V_pfr_second = F_A0_lev * np.trapz(inv_r_A_second[:idx_design], X_vals[:idx_design])

# CSTR volumes (rectangle)
V_cstr_first = F_A0_lev * X_design * inv_r_A_first[idx_design]
V_cstr_second = F_A0_lev * X_design * inv_r_A_second[idx_design]

fig, (ax1, ax2) = plt.subplots(1, 2, figsize=(12, 5))

# Levenspiel plot
ax1.plot(X_vals, inv_r_A_zero, 'g-', linewidth=2, label='Zero-order')
ax1.plot(X_vals, inv_r_A_first, 'b-', linewidth=2, label='First-order')
ax1.plot(X_vals, inv_r_A_second, 'r-', linewidth=2, label='Second-order')
ax1.axvline(x=X_design, color='gray', linestyle=':', alpha=0.5)
ax1.set_xlabel('Conversion $X_A$', fontsize=11)
ax1.set_ylabel('$-1/r_A$ (min·L/mol)', fontsize=11)
ax1.set_title('Levenspiel Plot for Reactor Comparison', fontsize=12)
ax1.legend(fontsize=10)
ax1.grid(True, alpha=0.3)
ax1.set_ylim(0, 80)

# Volume comparison
reactor_types = ['PFR\n(1st)', 'CSTR\n(1st)', 'PFR\n(2nd)', 'CSTR\n(2nd)']
volumes = [V_pfr_first, V_cstr_first, V_pfr_second, V_cstr_second]
colors = ['blue', 'lightblue', 'red', 'lightcoral']

bars = ax2.bar(reactor_types, volumes, color=colors, alpha=0.7, edgecolor='black')
ax2.set_ylabel('Required Volume (L)', fontsize=11)
ax2.set_title(f'Reactor Volume for $X_A = {X_design}$', fontsize=12)
ax2.grid(True, alpha=0.3, axis='y')

# Add value labels on bars
for bar, vol in zip(bars, volumes):
    height = bar.get_height()
    ax2.text(bar.get_x() + bar.get_width()/2., height,
            f'{vol:.1f}', ha='center', va='bottom', fontsize=10)

plt.tight_layout()
plt.savefig('reaction_engineering_plot5.pdf', dpi=150, bbox_inches='tight')
plt.close()

# Calculate efficiency ratios
pfr_cstr_ratio_first = V_cstr_first / V_pfr_first
pfr_cstr_ratio_second = V_cstr_second / V_pfr_second
\end{pycode}

\begin{figure}[H]
\centering
\includegraphics[width=0.95\textwidth]{reaction_engineering_plot5.pdf}
\caption{Levenspiel plot demonstrating graphical reactor design method. The area under each curve represents PFR volume, while CSTR volume corresponds to the rectangular area. For first-order kinetics at $X_A = \py{X_design}$, CSTR requires $\py{f'{pfr_cstr_ratio_first:.2f}'}$× more volume than PFR. Second-order reactions show even greater disparity ($\py{f'{pfr_cstr_ratio_second:.2f}'}$×), while zero-order reactions perform identically in both configurations.}
\end{figure}

\section{Reactor Performance Optimization}

The optimal reactor configuration depends on multiple factors:
\begin{itemize}
\item \textbf{Reaction order}: PFR advantage increases with positive reaction order
\item \textbf{Conversion target}: Higher conversions favor PFR over CSTR
\item \textbf{Heat management}: Exothermic reactions may require staged cooling
\item \textbf{Mixing requirements}: Fast reactions benefit from CSTR back-mixing
\item \textbf{Capital vs operating costs}: Multiple small CSTRs vs single large PFR
\end{itemize}

\begin{pycode}
# Comprehensive reactor performance comparison
conversion_targets = np.array([0.50, 0.70, 0.80, 0.90, 0.95])
k_opt = 0.1  # rate constant (1/min)
Q_opt = 10.0  # flow rate (L/min)
C_A0_opt = 1.0  # inlet concentration (mol/L)
F_A0_opt = Q_opt * C_A0_opt

# Calculate required volumes for each configuration
V_pfr_opt = []
V_cstr_1_opt = []
V_cstr_3_opt = []

for X in conversion_targets:
    # PFR
    V_pfr = (F_A0_opt / k_opt) * np.log(1 / (1 - X))
    V_pfr_opt.append(V_pfr)

    # Single CSTR
    V_cstr = (F_A0_opt * X) / (k_opt * (1 - X))
    V_cstr_1_opt.append(V_cstr)

    # 3 CSTRs in series
    def eq_3cstr(tau):
        return 1 - 1/(1 + k_opt*tau/3)**3 - X
    tau_3 = brentq(eq_3cstr, 0, 100)
    V_cstr_3_opt.append(tau_3 * Q_opt)

V_pfr_opt = np.array(V_pfr_opt)
V_cstr_1_opt = np.array(V_cstr_1_opt)
V_cstr_3_opt = np.array(V_cstr_3_opt)

fig, (ax1, ax2) = plt.subplots(1, 2, figsize=(12, 5))

# Volume comparison
ax1.plot(conversion_targets * 100, V_pfr_opt, 'b-o', linewidth=2, markersize=6, label='PFR')
ax1.plot(conversion_targets * 100, V_cstr_1_opt, 'r--s', linewidth=2, markersize=6, label='Single CSTR')
ax1.plot(conversion_targets * 100, V_cstr_3_opt, 'g-.^', linewidth=2, markersize=6, label='3 CSTRs in series')
ax1.set_xlabel('Conversion (\\%)', fontsize=11)
ax1.set_ylabel('Required Volume (L)', fontsize=11)
ax1.set_title('Reactor Volume vs Conversion', fontsize=12)
ax1.legend(fontsize=10)
ax1.grid(True, alpha=0.3)

# Efficiency ratios (CSTR/PFR)
ratio_cstr_1 = V_cstr_1_opt / V_pfr_opt
ratio_cstr_3 = V_cstr_3_opt / V_pfr_opt

ax2.plot(conversion_targets * 100, ratio_cstr_1, 'r--s', linewidth=2, markersize=6, label='Single CSTR / PFR')
ax2.plot(conversion_targets * 100, ratio_cstr_3, 'g-.^', linewidth=2, markersize=6, label='3 CSTRs / PFR')
ax2.axhline(y=1.0, color='b', linestyle='-', alpha=0.5, label='PFR baseline')
ax2.set_xlabel('Conversion (\\%)', fontsize=11)
ax2.set_ylabel('Volume Ratio (relative to PFR)', fontsize=11)
ax2.set_title('Reactor Efficiency Comparison', fontsize=12)
ax2.legend(fontsize=10)
ax2.grid(True, alpha=0.3)

plt.tight_layout()
plt.savefig('reaction_engineering_plot6.pdf', dpi=150, bbox_inches='tight')
plt.close()

# Store key optimization metrics
V_90_pfr = V_pfr_opt[-2]
V_90_cstr1 = V_cstr_1_opt[-2]
V_90_cstr3 = V_cstr_3_opt[-2]
savings_3cstr = (V_90_cstr1 - V_90_cstr3) / V_90_cstr1 * 100
\end{pycode}

\begin{figure}[H]
\centering
\includegraphics[width=0.95\textwidth]{reaction_engineering_plot6.pdf}
\caption{Reactor optimization analysis across conversion targets. The efficiency gap between CSTR and PFR widens dramatically at high conversions - at 90\% conversion, single CSTR requires $\py{f'{V_90_cstr1:.1f}'}$ L versus $\py{f'{V_90_pfr:.1f}'}$ L for PFR. Staging with 3 CSTRs in series reduces volume to $\py{f'{V_90_cstr3:.1f}'}$ L, achieving $\py{f'{savings_3cstr:.1f}'}$\% savings while preserving mixing benefits for heat management and fast reactions.}
\end{figure}

\section{Results Summary}

\begin{pycode}
# Compile key results from all analyses
results_data = [
    ['Batch (1st-order)', f'{t_90_first:.1f}', 'min', 'Time for 90\\% conversion'],
    ['Batch (2nd-order)', f'{t_90_second:.1f}', 'min', 'Time for 90\\% conversion'],
    ['CSTR (single)', f'{tau_single:.1f}', 'min', 'Residence time for 90\\%'],
    ['CSTR staging reduction', f'{reduction_pct:.1f}', '\\%', '3 stages vs single'],
    ['PFR volume (90\\%)', f'{V_pfr_90:.1f}', 'L', 'First-order reaction'],
    ['CSTR/PFR ratio', f'{volume_ratio:.1f}', '--', 'At 90\\% conversion'],
    ['Arrhenius $E_a$', f'{Ea_arr/1000:.0f}', 'kJ/mol', 'Activation energy'],
    ['$k$ ratio (25-100°C)', f'{k_ratio:.1f}', '--', 'Rate constant increase'],
    ['Time savings (25-50°C)', f'{time_reduction:.1f}', '\\%', 'Batch operation'],
]

print(r'\begin{table}[H]')
print(r'\centering')
print(r'\caption{Summary of Key Reactor Design Results}')
print(r'\begin{tabular}{@{}llll@{}}')
print(r'\toprule')
print(r'Parameter & Value & Units & Description \\')
print(r'\midrule')
for row in results_data:
    print(f"{row[0]} & {row[1]} & {row[2]} & {row[3]} \\\\")
print(r'\bottomrule')
print(r'\end{tabular}')
print(r'\end{table}')
\end{pycode}

\section{Conclusions}

This comprehensive computational analysis of chemical reactor design demonstrates several fundamental principles:

\textbf{Reactor Configuration}: For first-order irreversible reactions at 90\% conversion, PFR achieves residence time of $\tau = \py{f'{V_pfr_90/Q_pfr:.1f}'}$ min, while single CSTR requires $\tau = \py{f'{tau_single:.1f}'}$ min - a $\py{f'{volume_ratio:.1f}'}$× difference. CSTR staging (3 reactors in series) reduces this gap by $\py{f'{reduction_pct:.1f}'}$\%, approaching PFR efficiency while maintaining mixing advantages.

\textbf{Reaction Order Effects}: Levenspiel analysis reveals that higher-order reactions show greater disparity between reactor types. Second-order reactions require $\py{f'{pfr_cstr_ratio_second:.2f}'}$× more CSTR volume than PFR at $X_A = \py{X_design}$, while zero-order kinetics show no preference.

\textbf{Temperature Management}: Arrhenius behavior ($E_a = \py{f'{Ea_arr/1000:.0f}'}$ kJ/mol) demonstrates that increasing temperature from 25°C to 50°C reduces batch time by $\py{f'{time_reduction:.1f}'}$\%. Non-isothermal PFR analysis shows temperature rise of $\py{f'{T_rise:.1f}'}$ K for exothermic reactions, requiring active cooling to prevent thermal runaway and maintain selectivity.

\textbf{Design Optimization}: The optimal configuration balances capital costs (reactor volume), operating costs (temperature control), and process requirements (conversion, selectivity). Staged CSTR systems offer a practical compromise, achieving $\py{f'{savings_3cstr:.1f}'}$\% volume reduction versus single CSTR while preserving thermal management flexibility for fast exothermic reactions.

These numerical methods provide essential tools for reactor scale-up, process intensification, and economic optimization in chemical manufacturing \cite{Fogler2016, Schmidt2005}.

\begin{thebibliography}{99}

\bibitem{Fogler2016}
Fogler, H. S. (2016). \textit{Elements of Chemical Reaction Engineering} (5th ed.). Prentice Hall.

\bibitem{Levenspiel1999}
Levenspiel, O. (1999). \textit{Chemical Reaction Engineering} (3rd ed.). John Wiley \& Sons.

\bibitem{Smith2005}
Smith, J. M., Van Ness, H. C., \& Abbott, M. M. (2005). \textit{Introduction to Chemical Engineering Thermodynamics} (7th ed.). McGraw-Hill.

\bibitem{Scott2006}
Scott, D. M. (2006). \textit{Industrial Process Systems}. Springer.

\bibitem{Schmidt2005}
Schmidt, L. D. (2005). \textit{The Engineering of Chemical Reactions} (2nd ed.). Oxford University Press.

\bibitem{Aris1989}
Aris, R. (1989). \textit{Elementary Chemical Reactor Analysis}. Butterworth-Heinemann.

\bibitem{Davis2003}
Davis, M. E., \& Davis, R. J. (2003). \textit{Fundamentals of Chemical Reaction Engineering}. McGraw-Hill.

\bibitem{Froment2011}
Froment, G. F., Bischoff, K. B., \& De Wilde, J. (2011). \textit{Chemical Reactor Analysis and Design} (3rd ed.). John Wiley \& Sons.

\bibitem{Carberry2001}
Carberry, J. J. (2001). \textit{Chemical and Catalytic Reaction Engineering}. Dover Publications.

\bibitem{Hill1977}
Hill, C. G. (1977). \textit{An Introduction to Chemical Engineering Kinetics and Reactor Design}. John Wiley \& Sons.

\bibitem{Walas1991}
Walas, S. M. (1991). \textit{Modeling with Differential Equations in Chemical Engineering}. Butterworth-Heinemann.

\bibitem{Denbigh1981}
Denbigh, K. G., \& Turner, J. C. R. (1981). \textit{Chemical Reactor Theory: An Introduction} (3rd ed.). Cambridge University Press.

\bibitem{Holland1975}
Holland, C. D., \& Anthony, R. G. (1975). \textit{Fundamentals of Chemical Reaction Engineering}. Prentice Hall.

\bibitem{Westerterp1984}
Westerterp, K. R., Van Swaaij, W. P. M., \& Beenackers, A. A. C. M. (1984). \textit{Chemical Reactor Design and Operation} (2nd ed.). John Wiley \& Sons.

\bibitem{Nauman2008}
Nauman, E. B. (2008). \textit{Chemical Reactor Design, Optimization, and Scaleup} (2nd ed.). John Wiley \& Sons.

\bibitem{Rawlings2002}
Rawlings, J. B., \& Ekerdt, J. G. (2002). \textit{Chemical Reactor Analysis and Design Fundamentals}. Nob Hill Publishing.

\bibitem{Missen1999}
Missen, R. W., Mims, C. A., \& Saville, B. A. (1999). \textit{Introduction to Chemical Reaction Engineering and Kinetics}. John Wiley \& Sons.

\bibitem{Butt2000}
Butt, J. B. (2000). \textit{Reaction Kinetics and Reactor Design} (2nd ed.). Marcel Dekker.

\bibitem{Trambouze2004}
Trambouze, P., Van Landeghem, H., \& Wauquier, J. P. (2004). \textit{Chemical Reactors: Design, Engineering, Operation}. Editions Technip.

\bibitem{Ancheyta2017}
Ancheyta, J. (2017). \textit{Chemical Reaction Kinetics: Concepts, Methods and Case Studies}. John Wiley \& Sons.

\end{thebibliography}

\end{document}
